% 1-37 题 图1_32：惠更斯等时摆——上滚轮线和下滚轮线
% 轮半径 R=1，上滚轮线 P 点轨迹；下滚轮线关于 MN 翻转
\begin{tikzpicture}[>=Stealth, thick]
    % 中心水平线 MN
    \draw (-2.5, 0) -- (7, 0);
    \node[left] at (-2.5, 0) {$M$};
    \node[right] at (7, 0) {$N$};

    % 上下两个圆 (半径设为 1.2)
    \coordinate (O1) at (0, 1.2);
    \coordinate (O2) at (0, -1.2);
    \draw (O1) circle (1.2);
    \draw (O2) circle (1.2);

    % 中心点 P
    \node[below left] at (0, 0) {$P$};
    \node[above left] at (0, 0) {$P$}; % 原图似乎上下都标了 P 或有点模糊，放在左侧居中

    % 绘制椭圆 (中心偏右，穿过两圆)
    \draw (2.5, 0) ellipse (4.2 and 1.8);

    % 虚线水平切线
    \draw[dashed] (0, 2.4) -- (3, 2.4);
    \draw[dashed] (0, -2.4) -- (3, -2.4);

    % 圆的半径箭头 R
    \draw[->] (O1) -- (-0.85, 2.05) node[midway, right, xshift=-2pt, yshift=2pt] {$R$};
    \draw[->] (O2) -- (-0.85, -2.05) node[midway, right, xshift=-2pt, yshift=-2pt] {$R$};

    % 旋转方向的弧形箭头
    \draw[->] (0.8, 2.6) arc[start angle=60, end angle=120, radius=1.6];
    \draw[->] (-0.8, -2.6) arc[start angle=240, end angle=300, radius=1.6];
\end{tikzpicture}
